% Use a standard class when IEEEtran is not available
\documentclass[a4paper,10pt]{article}
\usepackage{cite}
\usepackage{amsmath,amssymb}
\usepackage{graphicx}
\usepackage{booktabs}
\usepackage{array}
\usepackage{url}

\begin{document}

\title{Voice-First LoRa Mesh Networking: Adaptive Spreading Factor and
Bandwidth Control with Voice-Aware Routing for Off-Grid
Emergency Communication}

\author{
\IEEEauthorblockN{Nadeesha N. J.}
\IEEEauthorblockA{Department of Electronic and Telecommunication Engineering\\
University of Moratuwa\\
Katubedda, Sri Lanka\\
Email: your.email@uom.lk}
}

\maketitle

\begin{abstract}
% TODO: 150-200 words. Write last.
Placeholder abstract.
\end{abstract}

\begin{IEEEkeywords}
LoRa, mesh networking, Codec2, adaptive modulation, emergency communication
\end{IEEEkeywords}

\section{Introduction}
% TODO
Placeholder.

\section{Literature Review}
% TODO
Placeholder.

\section{Materials and Methods}
% TODO
Placeholder.

%=====================================================================
%  RESULTS AND DISCUSSION  (submitted as plan / expected approach)
%=====================================================================
\section{Results and Discussion}

At the time of writing, the custom hardware platform is in the final
stage of schematic development and measurement data is not yet
available. This section therefore presents the evaluation plan: the
metrics that will be measured, the experimental procedure used to
obtain them, the baselines against which the system will be compared,
and the outcomes that are expected from each experiment.

\subsection{Evaluation Metrics}

The central claim of this work is that a LoRa link can be tuned so
that it delivers the highest possible speech quality for the airtime
that it consumes. To make this claim measurable, a composite metric,
the \emph{voice quality per unit airtime} $\eta$, is defined as

\begin{equation}
\eta = \frac{Q}{T_{\mathrm{air}}},
\end{equation}

\noindent where $Q$ is the objective speech quality score of the
received audio on a five point scale and $T_{\mathrm{air}}$ is the
measured on-air time in seconds required to transmit one second of
speech. A value of $\eta$ greater than the value obtained for a fixed
configuration indicates that the adaptive scheme has recovered
quality without spending extra airtime.

The following quantities will be logged for every trial.

\begin{itemize}
\item \textbf{Speech quality $Q$.} Objective scoring of the decoded
      audio against the transmitted reference using ViSQOL, supported
      by a subjective mean opinion score (MOS) panel of at least
      fifteen listeners for a representative subset of conditions.
\item \textbf{Intelligibility.} Word recognition rate obtained from a
      modified rhyme test, since a compressed speech codec can score
      poorly on quality yet remain fully understandable.
\item \textbf{Airtime $T_{\mathrm{air}}$.} Computed from the LoRa
      time-on-air expression for the selected spreading factor,
      bandwidth, coding rate and payload length, and cross checked
      against a GPIO toggle captured on a logic analyser at the
      transmit enable pin.
\item \textbf{Mouth to ear latency.} Measured end to end with a
      timestamped audio marker, reported as the median and the
      ninety-fifth percentile.
\item \textbf{Packet delivery ratio (PDR)} and \textbf{frame erasure
      rate}, together with the RSSI and SNR reported by the
      transceiver for each received frame.
\item \textbf{Transmissions per delivered frame}, which captures the
      cost of redundancy in the mesh layer.
\item \textbf{Channel occupancy}, the fraction of wall clock time the
      shared channel is busy across all nodes.
\end{itemize}

\subsection{Experiment 1: Spreading Factor and Bandwidth Sweep}

A two node point to point link will be characterised first, because
the mesh results cannot be interpreted without a reliable model of the
single hop behaviour. One node transmits a fixed set of speech samples
drawn from a standard corpus, encoded at each of the available Codec2
bit rates. The receiving node decodes and stores the audio for offline
scoring.

The parameter space is swept exhaustively: spreading factors SF5
through SF12, bandwidths of 203, 406, 812 and 1625~kHz, coding rates
of 4/5 through 4/8, and Codec2 modes at 700, 1200, 1600 and
3200~bit/s. Each combination is repeated thirty times at each of five
link distances, chosen to span the range from a strong link to the
edge of the coverage area. Transmit power is held constant so that
only modulation parameters vary. Results are reported as the mean with
a ninety-five percent confidence interval.

The outcome of this experiment is a lookup surface of $\eta$ over the
parameter space, indexed by the measured SNR. This surface becomes the
decision table used by the adaptive controller.

\begin{table}[t]
\caption{Planned test matrix for the link characterisation sweep}
\label{tab:matrix}
\centering
\renewcommand{\arraystretch}{1.2}
\begin{tabular}{p{0.30\columnwidth} p{0.55\columnwidth}}
\toprule
\textbf{Parameter} & \textbf{Levels} \\
\midrule
Spreading factor & SF5 to SF12 \\
Bandwidth        & 203, 406, 812, 1625~kHz \\
Coding rate      & 4/5, 4/6, 4/7, 4/8 \\
Codec2 mode      & 700, 1200, 1600, 3200~bit/s \\
Link distance    & 5 levels, strong link to link edge \\
Environment      & Open field, built-up campus, indoor \\
Repetitions      & 30 per cell \\
\bottomrule
\end{tabular}
\end{table}

\subsection{Experiment 2: Closed Loop Adaptive Control}

The adaptive controller selects the spreading factor, bandwidth and
codec rate from the surface produced in Experiment~1, using the SNR
and frame erasure rate of recent receptions as its input. Its
behaviour will be evaluated on a walking test in which one node moves
away from a stationary node at a constant pace until the link fails,
then returns. Three schemes are compared under the same route and the
same speech material:

\begin{enumerate}
\item a fixed conservative configuration, high spreading factor and
      narrow bandwidth, representing a link tuned for worst case
      range;
\item a fixed aggressive configuration, low spreading factor and wide
      bandwidth, representing a link tuned for quality;
\item the proposed adaptive controller.
\end{enumerate}

The expected result is that the adaptive scheme matches the aggressive
configuration for $\eta$ while the link is strong, and matches the
conservative configuration for PDR once the link degrades, so that the
usable communication range is extended without a proportional loss of
speech quality. The cost that must be reported honestly is the
settling time of the controller, that is, how many seconds of degraded
or lost audio occur before the correct configuration is reached after
a sudden change in channel conditions.

\subsection{Experiment 3: Voice-Aware Mesh Routing}

A testbed of six to eight nodes will be deployed so that the network
diameter is at least three hops. Existing LoRa mesh implementations
are text oriented and treat every packet identically, so a managed
flooding scheme of the kind used by Meshtastic is adopted as the
baseline. The proposed protocol differs in that it classifies traffic,
gives voice frames a short deadline and drops them rather than queuing
them once that deadline has passed, and suppresses rebroadcast of a
voice frame by a neighbour that is not on a path toward the
destination.

Three traffic loads will be applied: voice only, voice mixed with
periodic position reports, and a congested load in which two voice
streams share the network. For each load the protocols are compared on
delivery ratio, ninety-fifth percentile latency, transmissions per
delivered frame and channel occupancy.

The expected result is a reduction in transmissions per delivered
frame and in tail latency relative to managed flooding, obtained
because stale voice frames are discarded instead of consuming airtime
that a fresher frame needs. The anticipated trade-off is a lower
delivery ratio for non-urgent text traffic during voice activity,
which will be quantified rather than avoided.

\subsection{Threats to Validity}

Three risks are recognised. First, objective speech quality metrics
are validated mainly on telephone band codecs and may correlate poorly
with listener judgement at very low bit rates, which is why the
subjective panel is retained. Second, outdoor field results depend on
weather, foliage and interference in the shared 2.4~GHz band, so every
comparison is run within a single session and interference is logged
with a spectrum capture. Third, a testbed of eight nodes cannot
demonstrate behaviour at the scale of a large deployment, so the
measured per hop behaviour will be used to parameterise a simulation
for the larger case, and the simulation will be reported as an
extrapolation rather than as a measurement.

\section{Conclusion}
% TODO
Placeholder.

\section*{Acknowledgment}
% TODO

\begin{thebibliography}{1}
\bibitem{ref1}
Placeholder reference.
\end{thebibliography}

\end{document}